\documentclass[11pt]{article}
\usepackage{fontspec,amsmath,amssymb,amsthm,hyperref}
\usepackage[margin=1in]{geometry}
\setlength{\emergencystretch}{3em}
\title{Random-center smoothing and the integrated Gaussian energy bound}
\author{LangevinLean mathematical exposition}
\date{}
\begin{document}\maketitle
\noindent Source module: \texttt{AverageLMC/GaussianPerturbation.lean}. Compile this proof document with XeLaTeX or LuaLaTeX. The accompanying HTML pairs the mathematical argument with the exact Lean statements and records the scope of the formalization.


This module carries out the expectation and time-integration steps used in the improved one-step estimate. Independence promotes a deterministic smoothing estimate to a random center; strong convexity and orthogonality then give the exact Gaussian energy constant.



The energy theorem uses an affine upper bound for the expected potential difference, the exact Brownian increment moment, orthogonality, and a common lower bound \(w\) for \(\|X+B_h-Y_t\|_{L^2}^2\). Its key input and conclusion are


\[\mathbb Ef(X+B_t)-\mathbb Ef(Y_t)\le a+bt\qquad(0\le t\le h).\]
\[(1+mh)w\le\|X-Y_0\|_{L^2}^{2}+2ha+h^2b+mh^2d.\]
\section{Promote smoothing to an independent random center}


For independent L\ensuremath{{}^{2}} vectors X and Z, their joint pushforward measure is the product of their marginal laws. The function F(x,z)=f(x+c z)−f(x) is continuous and integrable under this joint law because both arguments have finite second moment and f has quadratic growth. The pushforward identity and Fubini turn its expectation into the iterated integral. Integrate the pointwise smoothing bound over the law of X, which is a probability measure, to retain the same constant. Finally use almost-everywhere representatives to identify the pointwise sum with the sum in L\ensuremath{{}^{2}}.



\noindent Lean declarations: \texttt{independent\_smoothing}.

\section{Integrate strong convexity with the future-noise correction}


Let w be a common lower bound for the squared coupling costs between X+B(h) and the stationary comparison Y(t). Orthogonality expands the latter cost exactly as ‖X+B(t)−Y(t)‖\ensuremath{{}^{2}}+2(h−t)d. Integrate the strong supporting-plane bound, using an affine upper bound a+b t for the potential expectation difference. The integrand is bounded by 2a+2mhd−mw+(2b−2md)t. Its explicit integral is 2ha+mh\ensuremath{{}^{2}}d−mhw+h\ensuremath{{}^{2}}b. Combine this with the exact synchronous energy identity and the endpoint coupling bound to obtain (1+mh)w\ensuremath{\le}‖X−Y(0)‖\ensuremath{{}^{2}}+2ha+h\ensuremath{{}^{2}}b+mh\ensuremath{{}^{2}}d. Every input involving noise or stationarity is discharged for the actual Gaussian step by the canonical construction in the following modules.



\noindent Lean declarations: \texttt{gaussian\_energy\_bound}.

\end{document}
